\documentclass[11pt,a4paper]{article}
% Wspólny szablon dla kart supermocy (jedna karta = jedna strona A4)
% Użycie w karcie: \documentclass[11pt,a4paper]{article} % Wspólny szablon dla kart supermocy (jedna karta = jedna strona A4)
% Użycie w karcie: \documentclass[11pt,a4paper]{article} % Wspólny szablon dla kart supermocy (jedna karta = jedna strona A4)
% Użycie w karcie: \documentclass[11pt,a4paper]{article} \input{../preambula} \begin{document}

% Silnik tectonic (XeTeX): polskie znaki przez fontspec, nie przez inputenc/T1
\usepackage{fontspec}
\setmainfont{Arial}
\usepackage[polish]{babel}
\usepackage[a4paper,margin=1.5cm]{geometry}
\usepackage{amsmath,amssymb}
\usepackage{tikz}
\usetikzlibrary{arrows.meta,positioning,calc,shapes.geometric}
\usepackage[most]{tcolorbox}
\usepackage{enumitem}
\usepackage{booktabs}
\usepackage{xcolor}

\pagestyle{empty}
\setlength{\parindent}{0pt}
\setlength{\parskip}{0.4em}

\definecolor{cel}{HTML}{1F4E79}
\definecolor{regula}{HTML}{E8F1FA}
\definecolor{pulapka}{HTML}{FDECEA}
\definecolor{okres}{HTML}{C0392B}
\definecolor{skonczony}{HTML}{1E8449}

% Ramka na regułę
\newtcolorbox{regula}[1][]{
  colback=regula, colframe=cel, boxrule=0.8pt, arc=2pt,
  left=6pt, right=6pt, top=4pt, bottom=4pt, title=#1,
  fonttitle=\bfseries, coltitle=white
}

% Ramka na pułapkę z kartkówki
\newtcolorbox{pulapka}[1][]{
  colback=pulapka, colframe=okres, boxrule=0.8pt, arc=2pt,
  left=6pt, right=6pt, top=4pt, bottom=4pt, title=#1,
  fonttitle=\bfseries, coltitle=white
}

% Nagłówek karty: numer, tytuł, cel
\newcommand{\naglowekKarty}[3]{%
  {\LARGE\bfseries\color{cel} Karta #1: #2}\\[2pt]
  {\large Cel: #3}\par
  \vspace{0.3em}\noindent\rule{\linewidth}{1pt}\par
}

% Pole na obliczenia ucznia: ramka z delikatną kratką, argument = wysokość
\newcommand{\miejsce}[1]{%
  \par\vspace{0.2em}\noindent
  \begin{tikzpicture}
    \draw[step=0.5cm, gray!25, very thin] (0,0) grid (\textwidth,#1);
    \draw[gray!60, rounded corners=2pt] (0,0) rectangle (\textwidth,#1);
    \node[anchor=north west, font=\scriptsize, text=gray] at (0.05,#1) {obliczenia};
  \end{tikzpicture}\par}
 \begin{document}

% Silnik tectonic (XeTeX): polskie znaki przez fontspec, nie przez inputenc/T1
\usepackage{fontspec}
\setmainfont{Arial}
\usepackage[polish]{babel}
\usepackage[a4paper,margin=1.5cm]{geometry}
\usepackage{amsmath,amssymb}
\usepackage{tikz}
\usetikzlibrary{arrows.meta,positioning,calc,shapes.geometric}
\usepackage[most]{tcolorbox}
\usepackage{enumitem}
\usepackage{booktabs}
\usepackage{xcolor}

\pagestyle{empty}
\setlength{\parindent}{0pt}
\setlength{\parskip}{0.4em}

\definecolor{cel}{HTML}{1F4E79}
\definecolor{regula}{HTML}{E8F1FA}
\definecolor{pulapka}{HTML}{FDECEA}
\definecolor{okres}{HTML}{C0392B}
\definecolor{skonczony}{HTML}{1E8449}

% Ramka na regułę
\newtcolorbox{regula}[1][]{
  colback=regula, colframe=cel, boxrule=0.8pt, arc=2pt,
  left=6pt, right=6pt, top=4pt, bottom=4pt, title=#1,
  fonttitle=\bfseries, coltitle=white
}

% Ramka na pułapkę z kartkówki
\newtcolorbox{pulapka}[1][]{
  colback=pulapka, colframe=okres, boxrule=0.8pt, arc=2pt,
  left=6pt, right=6pt, top=4pt, bottom=4pt, title=#1,
  fonttitle=\bfseries, coltitle=white
}

% Nagłówek karty: numer, tytuł, cel
\newcommand{\naglowekKarty}[3]{%
  {\LARGE\bfseries\color{cel} Karta #1: #2}\\[2pt]
  {\large Cel: #3}\par
  \vspace{0.3em}\noindent\rule{\linewidth}{1pt}\par
}

% Pole na obliczenia ucznia: ramka z delikatną kratką, argument = wysokość
\newcommand{\miejsce}[1]{%
  \par\vspace{0.2em}\noindent
  \begin{tikzpicture}
    \draw[step=0.5cm, gray!25, very thin] (0,0) grid (\textwidth,#1);
    \draw[gray!60, rounded corners=2pt] (0,0) rectangle (\textwidth,#1);
    \node[anchor=north west, font=\scriptsize, text=gray] at (0.05,#1) {obliczenia};
  \end{tikzpicture}\par}
 \begin{document}

% Silnik tectonic (XeTeX): polskie znaki przez fontspec, nie przez inputenc/T1
\usepackage{fontspec}
\setmainfont{Arial}
\usepackage[polish]{babel}
\usepackage[a4paper,margin=1.5cm]{geometry}
\usepackage{amsmath,amssymb}
\usepackage{tikz}
\usetikzlibrary{arrows.meta,positioning,calc,shapes.geometric}
\usepackage[most]{tcolorbox}
\usepackage{enumitem}
\usepackage{booktabs}
\usepackage{xcolor}

\pagestyle{empty}
\setlength{\parindent}{0pt}
\setlength{\parskip}{0.4em}

\definecolor{cel}{HTML}{1F4E79}
\definecolor{regula}{HTML}{E8F1FA}
\definecolor{pulapka}{HTML}{FDECEA}
\definecolor{okres}{HTML}{C0392B}
\definecolor{skonczony}{HTML}{1E8449}

% Ramka na regułę
\newtcolorbox{regula}[1][]{
  colback=regula, colframe=cel, boxrule=0.8pt, arc=2pt,
  left=6pt, right=6pt, top=4pt, bottom=4pt, title=#1,
  fonttitle=\bfseries, coltitle=white
}

% Ramka na pułapkę z kartkówki
\newtcolorbox{pulapka}[1][]{
  colback=pulapka, colframe=okres, boxrule=0.8pt, arc=2pt,
  left=6pt, right=6pt, top=4pt, bottom=4pt, title=#1,
  fonttitle=\bfseries, coltitle=white
}

% Nagłówek karty: numer, tytuł, cel
\newcommand{\naglowekKarty}[3]{%
  {\LARGE\bfseries\color{cel} Karta #1: #2}\\[2pt]
  {\large Cel: #3}\par
  \vspace{0.3em}\noindent\rule{\linewidth}{1pt}\par
}

% Pole na obliczenia ucznia: ramka z delikatną kratką, argument = wysokość
\newcommand{\miejsce}[1]{%
  \par\vspace{0.2em}\noindent
  \begin{tikzpicture}
    \draw[step=0.5cm, gray!25, very thin] (0,0) grid (\textwidth,#1);
    \draw[gray!60, rounded corners=2pt] (0,0) rectangle (\textwidth,#1);
    \node[anchor=north west, font=\scriptsize, text=gray] at (0.05,#1) {obliczenia};
  \end{tikzpicture}\par}


\begin{document}

\naglowekKarty{3}{Porównywanie ułamków}{wpisać $>$, $<$ lub $=$ między dwiema liczbami wymiernymi}

\vspace{0.3em}
\begin{regula}[REGUŁA – WYBIERZ JEDNĄ Z CZTERECH DRÓG]
\begin{enumerate}[leftmargin=*, itemsep=2pt]
  \item \textbf{Dziesiętne:} wyrównaj liczbę cyfr po przecinku zerami. $0{,}6$ vs $0{,}57$: $0{,}60 > 0{,}57$.
  \item \textbf{Okres:} rozpisz okres, aż będzie widać różnicę. $0{,}(4)=0{,}444\ldots > 0{,}400 = 0{,}4$.
  \item \textbf{Wspólny mianownik} (ułamki zwykłe): rozszerz oba ułamki do tego samego mianownika i porównaj liczniki. $\dfrac{3}{4}$ vs $\dfrac{5}{7}$: $\dfrac{21}{28}>\dfrac{20}{28}$.
  \item \textbf{Ta sama góra / ten sam dół:} przy tym samym liczniku większy jest mniejszy mianownik ($\dfrac{3}{7}>\dfrac{3}{8}$).
\end{enumerate}
\end{regula}

\vspace{0.4em}
\textbf{Ułamki zwykłe krok po kroku: wspólny mianownik. Druga droga: oś liczbowa.} Porównajmy $\dfrac{5}{6}$ i $\dfrac{7}{9}$:
\begin{enumerate}[leftmargin=*, itemsep=3pt]
  \item \textbf{Wspólny mianownik:} wypisz wielokrotności $6$: $6, 12, \mathbf{18}$ oraz $9$: $9, \mathbf{18}$. Najmniejsza wspólna to $18$.
  \item \textbf{Rozszerz oba ułamki do mianownika $18$:} $\dfrac{5}{6}=\dfrac{5\cdot 3}{6\cdot 3}=\dfrac{15}{18}$ \quad oraz \quad $\dfrac{7}{9}=\dfrac{7\cdot 2}{9\cdot 2}=\dfrac{14}{18}$.
  \item \textbf{Porównaj liczniki} (mianowniki są takie same): $15>14$, więc $\dfrac{5}{6}>\dfrac{7}{9}$.
\end{enumerate}
Na osi liczbowej, po podzieleniu jej na osiemnaste części, widać to samo: większy licznik leży bardziej na prawo.

\begin{center}
\begin{tikzpicture}[>=Stealth, x=1.6cm]
  \draw[->] (-0.3,0) -- (6.5,0) node[right, font=\small] {$x$};
  \foreach \k in {12,...,18} {
    \pgfmathsetmacro{\px}{\k-12}
    \draw (\px,0.1) -- (\px,-0.1);
    \node[below, font=\small] at (\px,-0.1) {$\frac{\k}{18}$};
  }
  \fill[cel] (2,0) circle (2.2pt) node[above left=2pt, font=\small\bfseries, color=cel] {$\frac{14}{18}=\frac{7}{9}$};
  \fill[okres] (3,0) circle (2.2pt) node[above right=2pt, font=\small\bfseries, color=okres] {$\frac{15}{18}=\frac{5}{6}$};
  \node[font=\small, color=okres] at (3,-1.0) {$\frac{5}{6}$ leży bardziej na prawo};
\end{tikzpicture}
\end{center}

\vspace{0.2em}
\textbf{Liczby ujemne (termometr):} im bliżej zera, tym większa liczba. $-0{,}2$ jest bliżej zera niż $-0{,}25$, więc $-0{,}2>-0{,}25$.

\vspace{0.2em}
\textbf{Przykłady rozstrzygnięte:}
\begin{center}
\renewcommand{\arraystretch}{1.6}
\begin{tabular}{@{}lll@{}}
\toprule
Porównanie & Metoda & Wynik \\
\midrule
$0{,}6$ vs $0{,}57$ & wyrównanie zerami: $0{,}60$ vs $0{,}57$ & $0{,}6>0{,}57$ \\
$0{,}(4)$ vs $0{,}4$ & rozpisz okres: $0{,}444\ldots$ vs $0{,}400$ & $0{,}(4)>0{,}4$ \\
$\dfrac{5}{6}$ vs $\dfrac{7}{9}$ & mianownik $18$: $\dfrac{15}{18}$ vs $\dfrac{14}{18}$ & $\dfrac{5}{6}>\dfrac{7}{9}$ \\
$\dfrac{3}{4}$ vs $\dfrac{5}{7}$ & mianownik $28$: $\dfrac{21}{28}$ vs $\dfrac{20}{28}$ & $\dfrac{3}{4}>\dfrac{5}{7}$ \\
$\dfrac{3}{7}$ vs $\dfrac{3}{8}$ & ten sam licznik & $\dfrac{3}{7}>\dfrac{3}{8}$ \\
$-0{,}2$ vs $-0{,}25$ & termometr: bliżej zera & $-0{,}2>-0{,}25$ \\
\bottomrule
\end{tabular}
\end{center}

\begin{pulapka}[PUŁAPKI Z KARTKÓWKI]
\begin{enumerate}[leftmargin=*, itemsep=2pt]
  \item $0{,}4\neq 0{,}(4)$. Okres to nieskończony ciąg czwórek, więc $0{,}(4)$ jest większe.
  \item Więcej cyfr nie znaczy większa liczba: $0{,}7>0{,}68$, bo po wyrównaniu zerami $0{,}70>0{,}68$.
  \item Przy liczbach ujemnych odwróć intuicję: większa wartość bezwzględna oznacza \emph{mniejszą} liczbę.
\end{enumerate}
\end{pulapka}

\end{document}
